% chapters/07/sub-cpmk.tex
\begin{subcpmkbox}
Bab ini menuntaskan \SUBCPMK{1.4} secara utuh, sekaligus menutup \CPMK{1}.
Ketiga indikatornya diperiksa, dan buktinya berupa prototipe permainan yang
bendanya bertabrakan dengan benar.

\vspace{0.5em}
\textbf{\SUBCPMK{1.4} Mendeteksi dan menangani \istilah{collision} menggunakan
\code{Rect} dan/atau \istilah{mask} pada \istilah{sprite}} --- ranah C4,
\CPMK{1}, \CPL{03} dan \CPL{04}.

\vspace{0.5em}
{\setlength{\tabcolsep}{4pt}\footnotesize%
\begin{tabularx}{\textwidth}{@{}C{1.9cm} L{4.0cm} L{2.5cm} Y@{}}
\toprule
\headrow \textbf{Indikator} & \textbf{Rumusan} & \textbf{Dibahas pada} & \textbf{Bukti} \\
\midrule
\ind{a} & Memilih metode \istilah{collision} yang tepat sesuai kebutuhan &
Bagian~\ref{sec:deteksi-tabrakan} dan Bagian~\ref{sec:mask} &
Alasan pemilihan metode pada laporan praktikum, beserta penjelasan kapan
persegi pembatas tidak memadai. \\
\addlinespace
\ind{b} & Menyelesaikan tabrakan per sumbu sehingga objek tidak menembus dinding &
Bagian~\ref{sec:per-sumbu} &
Kode penyelesaian tabrakan per sumbu, beserta hasil pemeriksaan otomatis yang
menunjukkan tidak ada bingkai bertumpang tindih. \\
\addlinespace
\ind{c} & Menganalisis dan memperbaiki kasus tabrakan yang salah &
Bagian~\ref{sec:analisis-tabrakan} &
Analisis tiga gejala tabrakan yang salah beserta perbaikannya, dan catatan
perbaikan pada laporan praktikum. \\
\bottomrule
\end{tabularx}}
\end{subcpmkbox}
